% Agent 06: Book pages 365--368 (PDF pp.11--14)
% Section 2.5 Hydrodynamics and Thermodynamics
% Covers: basic references, notation, parameters describing the fluid,
%         fundamental thermodynamic relation, equation of state,
%         adiabatic index, speed of sound, first and second laws of
%         thermodynamics, decomposition of 4-velocity, frozen-in magnetic
%         field, stress-energy tensor, equations of hydrodynamics

%% ----------------------------------------------------------------
%% 2.5  Hydrodynamics and Thermodynamics
%% ----------------------------------------------------------------

\subsection{Hydrodynamics and Thermodynamics}
\label{sec:2.5}

\emph{Basic references}~~\S\S\,22.2 and 22.3 of Misner, Thorne, and Wheeler
(1973)~\cite{MTW1973}; Lichnerowicz (1967)~\cite{Lichnerowicz1967} and
references cited therein.  We adopt the notational conventions of
Misner, Thorne, and Wheeler (cited henceforth as MTW), including
signature ``$-\,+\,+\,+$''; geometrized units ($c = G = 1$); Greek
indices ranging from 0 to 3; and Latin from 1 to 3.  ``Hats'' on indices
refer to local orthonormal frames, e.g.\ $\mu^{\hat{\alpha}}$ and
$T^{\hat{\alpha}\hat{\beta}}$; extra-bold, sans-serif type for 4-vectors
and 4-tensors, e.g.\ $\mathbf{u}$, $\mathbf{p}^{\hat{\alpha}}$ and
$\mathbf{T}^{\hat{\alpha}\hat{\beta}}$; normal bold-face type for
3-vectors and 3-tensors (local Euclidean geometry).

\bigskip
\noindent
\emph{Parameters describing the ``fluid''}

\smallskip
\noindent
``LRF'' \quad local rest frame of the baryons; i.e.\ local orthonormal
frame in which there is no net baryon flux in any direction.

\smallskip
\noindent
$\mathbf{u}$ \quad 4-velocity of LRF; i.e.\ 4-velocity of the ``fluid''.

\smallskip
\noindent
$n$ \quad number density of baryons (number of baryons per unit volume)
as measured in LRF.

\smallskip
\noindent
$m_B$ \quad mean rest mass of a baryon in the fluid; i.e.
%
\begin{equation}
\tag{2.5.1}
m_B
= \left(\!\text{rest mass of}\atop\text{hydrogen atom}\!\right)
  \times \left(\!\text{fraction of all baryons that are}\atop
              \text{in the form of hydrogen nuclei}\!\right)
  + \frac{1}{4}\!\left(\!\text{rest mass of}\atop\text{helium atom}\!\right)
  \times \left(\!\text{fraction of all baryons that}\atop
              \text{are in helium nuclei}\!\right)
  + \cdots
\end{equation}

\emph{Note:} in this section we restrict ourselves to fluids that are
chemically homogeneous and in which no nuclear reactions occur
(``standard fluid'').  Thus, $m_B$ is constant.

\medskip
\noindent
$\rho_0$ \quad rest-mass density, defined by
%
\begin{equation}
\tag{2.5.2}
\rho_0 \equiv m_B\, n.
\end{equation}

\noindent
$\mathcal{V}$ \quad specific volume, i.e.\ volume per baryon, defined by
%
\begin{equation}
\tag{2.5.3}
\mathcal{V} = 1/n.
\end{equation}

\noindent
$V_0$ \quad specific volume,
%
\begin{equation}
\tag{2.5.4}
V_0 = 1/\rho_0.
\end{equation}

\noindent
$\rho$ \quad total density of mass-energy, as measured in LRF.

\noindent
$\Pi$ \quad specific internal energy, defined by
%
\begin{equation}
\tag{2.5.5}
\rho = \rho_0\,(1 + \Pi).
\end{equation}

\noindent
\emph{Note:} here and throughout this section we set the speed of light
equal to one.

\medskip
\noindent
$p$ \quad isotropic pressure, as measured in LRF.

\noindent
$T$ \quad temperature, as measured in LRF.

\noindent
$s$ \quad entropy per baryon, as measured in LRF; of course,

\noindent
$s_0$ \quad entropy per unit mass, as measured in LRF.
%
\begin{equation}
\tag{2.5.6}
s_0 = m_B^{-1}\, s.
\end{equation}

\noindent
$\mu$ \quad chemical potential, as measured in LRF; defined by
%
\begin{equation}
\tag{2.5.7}
\mu \equiv \left(\frac{\partial\rho}{\partial n}\right)_{\!s}
  = \frac{\rho + p}{n}\,.
\end{equation}

(The second equality follows from the first law of thermodynamics,
below.)

\medskip
\noindent
$\mathbf{q}$ \quad flux of energy (due to heat conduction, radiation,
convection, etc.)\ as measured in LRF.  This flux
(ergs\,cm$^{-2}$\,sec$^{-1}$) is a purely spatial vector as measured
in LRF; i.e.
%
\begin{equation}
\tag{2.5.8}
\mathbf{q} \cdot \mathbf{u} = 0.
\end{equation}

\noindent
$\mathbf{S}$ \quad entropy density-flux vector.  In LRF this vector has
time component equal to the entropy density, $S^0 = ns$, and space
components equal to the entropy flux, $S^j = q^j / T$.  Hence, in
frame-independent notation
%
\begin{equation}
\tag{2.5.9}
\mathbf{S} = ns\,\mathbf{u} + \mathbf{q}/T.
\end{equation}

\noindent
Note that
%
\begin{equation}
\tag{2.5.10}
\boldsymbol{\nabla}\cdot\mathbf{S}
= \left\{
\begin{array}{l}
\text{rate at which entropy is being generated}\\
\text{per unit volume as measured in LRF}
\end{array}
\right\}.
\end{equation}

\bigskip
\noindent
\emph{First law of thermodynamics}

\smallskip
\noindent
Follow a fluid element, containing $A$ baryons, along its world tube.
The total mass-energy in the fluid element, $\rho A/n$, changes as a
result of compression (change in volume, $A/n$) and as a result of
influx of heat:
%
\begin{equation}
\tag{2.5.11}
d(\rho A/n) = -p\,d(A/n) + T\,d(As).
\end{equation}

\noindent
[Here and throughout this section we assume for simplicity no internal
generation of entropy in the fluid element---e.g., no irreversible
chemical reactions; this enables us to write the influx of heat in
terms of the change in entropy, $T\,d(As)$.]

One often uses $A = \text{const.}$ and $\mu \equiv (\rho + p)/n$ to
rewrite this first law of thermodynamics in the equivalent forms
%
\begin{align}
\tag{2.5.12}
d\rho &= \frac{\rho + p}{n}\, dn + nT\, ds, \\[4pt]
\tag{2.5.12$'$}
d\mu  &= V\, dp + T\, ds.
\end{align}

\noindent
From the first law one can read off partial derivatives, e.g.
%
\begin{equation}
\tag{2.5.12$''$}
(\partial\rho/\partial n)_s = (\rho + p)/n, \qquad
(\partial\mu/\partial p)_s = T.
\end{equation}

\bigskip
\noindent
\emph{Fundamental relation and equations of state}

\smallskip
\noindent
The peculiar thermodynamic properties of the particular fluid being
studied are determined by a ``fundamental thermodynamic relation''
%
\begin{equation}
\tag{2.5.13}
\rho = \rho(n,\,s) \qquad \text{or} \qquad
\mu = \mu(p,\,s).
\end{equation}

\noindent
Once this relation has been specified---and once the constant $m_B$ has
been specified---one can use the first law of thermodynamics
(2.5.12) or (2.5.12$'$) to derive explicit expressions
for all other thermodynamic variables as functions of $n$, $s$ or $p$,
$s$.  For example, by combining with the first law one can derive the
``equations of state''
%
\begin{equation}
\tag{2.5.8$'$}
T(n,s) = \frac{1}{n}\left(\frac{\partial\rho}{\partial s}\right)_{\!n},
\qquad
p(n,s) = n\!\left(\frac{\partial\rho}{\partial n}\right)_{\!s} - \rho.
\end{equation}

\bigskip
\noindent
\emph{Adiabatic index and speed of sound}

\smallskip
\noindent
One defines the adiabatic index $\Gamma_1$ by
%
\begin{equation}
\tag{2.5.14}
\Gamma_1
\equiv \left(\frac{\partial\ln p}{\partial\ln n}\right)_{\!s}
= \frac{\rho + p}{p}\,
  \left(\frac{\partial p}{\partial\rho}\right)_{\!s}\,,
\end{equation}
%
where the third equality follows from the first law of thermodynamics.
It turns out that weak adiabatic perturbations (weak ``sound waves'')
propagate, in the LRF, with ordinary velocity
%
\begin{equation}
\tag{2.5.15}
c_s
= \left[\left(\frac{\partial p}{\partial\rho}\right)_{\!s}\right]^{1/2}
= \left(\frac{\Gamma_1\, p}{\rho + p}\right)^{\!1/2}.
\end{equation}

\bigskip
\noindent
\emph{Second law of thermodynamics}

\smallskip
\noindent
Equation~(2.5.10) allows one to write the second law of thermodynamics
in the form
%
\begin{equation}
\tag{2.5.16}
\boldsymbol{\nabla}\cdot\mathbf{S} \geq 0.
\end{equation}

\bigskip
\noindent
\emph{Decomposition of 4-velocity}

\smallskip
\noindent
One decomposes the gradient of the 4-velocity, $\nabla\mathbf{u}$,
into its ``irreducible'' tensorial parts:
%
\begin{equation}
\tag{2.5.17}
u_{\alpha;\beta}
= \sigma_{\alpha\beta} + \omega_{\alpha\beta}
  + \tfrac{1}{3}\,\theta\, h_{\alpha\beta}
  - a_\alpha\, u_\beta\,.
\end{equation}

\noindent
Here $\mathbf{a}$ is the 4-acceleration of the fluid
%
\begin{equation}
\tag{2.5.18a}
\mathbf{a} \equiv \nabla_{\mathbf{u}}\,\mathbf{u},
\qquad\text{i.e.}\quad
a^\alpha = u^\alpha{}_{;\beta}\, u^\beta\,,
\end{equation}
%
and $-a_\alpha u_\beta$ is that portion of $u_{\alpha;\beta}$ which is
not orthogonal to $\mathbf{u}$.  The remainder of $u_{\alpha;\beta}$
(i.e.\ $\sigma_{\alpha\beta} + \omega_{\alpha\beta} +
\tfrac{1}{3}\theta\, h_{\alpha\beta}$) is orthogonal to $\mathbf{u}$;
i.e.\ in the LRF it has only spatial components.  This orthogonal part
is decomposed into an isotropic part,
%
\begin{equation}
\tag{2.5.18b}
\theta \equiv \nabla \cdot \mathbf{u}
= u^\alpha{}_{;\alpha}\,,
\end{equation}
%
expansion, $\tfrac{1}{3}\theta h_{\alpha\beta}$, where $\theta$ is the
``expansion''; and $h_{\alpha\beta}$ is the ``projection tensor''
%
\begin{equation}
\tag{2.5.18c}
h_{\alpha\beta} \equiv g_{\alpha\beta} + u_\alpha\, u_\beta\,;
\end{equation}
%
plus a symmetric, trace-free ``shear''
%
\begin{equation}
\tag{2.5.18d}
\sigma_{\alpha\beta}
= \tfrac{1}{2}(u_{\alpha;\mu}\, h^\mu{}_\beta
              + u_{\beta;\mu}\, h^\mu{}_\alpha)
  - \tfrac{1}{3}\theta\, h_{\alpha\beta}\,;
\end{equation}
%
plus an antisymmetric ``rotation'' or ``vorticity''
%
\begin{equation}
\tag{2.5.18e}
\omega_{\alpha\beta}
= \tfrac{1}{2}(u_{\alpha;\mu}\, h^\mu{}_\beta
              - u_{\beta;\mu}\, h^\mu{}_\alpha).
\end{equation}

An observer in the LRF sees all fluid elements in his neighborhood to
move with low (nonrelativistic) velocities.  Let him use standard
Newtonian methods to calculate or measure the expansion $\theta$; shear
$\sigma^{\hat{j}}_{\;\hat{k}}$; and rotation $\omega^{\hat{j}}_{\;\hat{k}}$
of the fluid, and let a relativist calculate $\theta$,
$\sigma^{\hat{j}}_{\;\hat{k}}$, and
$\omega^{\hat{j}}_{\;\hat{k}}$ in the LRF from the above equations.
The two calculations will give the same answers.
[See, e.g., Ellis (1971).]

\bigskip
\noindent
\emph{Frozen-in magnetic field}

\smallskip
\noindent
Consider an ionized plasma which contains a ``frozen-in'' magnetic field.
The field is pure magnetic (no electric field) in the LRF.  Therefore
it can be described by a magnetic-field 4-vector $\mathbf{B}$
orthogonal to $\mathbf{u}$,
%
\begin{equation}
\tag{2.5.19}
\mathbf{B}\cdot\mathbf{u} = 0.
\end{equation}

\noindent
[For simplicity we assume that the magnetic permeability of the plasma
is the same as that of vacuum; so we do not distinguish between
``$\mathbf{B}$'' and ``$\mathbf{M}$''.  For a more general treatment, see
Lichnerowicz (1967).]  As the fluid moves, carrying with it the
frozen-in B-field, $\mathbf{B}$ must change as
%
\begin{equation}
\tag{2.5.20}
\frac{DB^\alpha}{d\tau}
= u_{\alpha;\beta}\, B^\beta
  + (\sigma_{\alpha\beta} - \tfrac{2}{3}\theta\, h_{\alpha\beta})\, B^\beta\,.
\end{equation}

\noindent
The term $u_{\alpha;\beta}\, B^\beta$ is required to keep $\mathbf{B}$
orthogonal to $\mathbf{u}$; by the term $\sigma_{\alpha\beta}\, B^\beta$,
the rotation of the fluid rotates the field lines; by the term
$-\tfrac{2}{3}\theta\, h_{\alpha\beta}\, B^\beta$, the compression of the
fluid orthogonal to $\mathbf{B}$ magnifies $\mathbf{B}$, conserving
flux.

\bigskip
\noindent
\emph{Stress-energy tensor}

\smallskip
\noindent
Consider a fluid with isotropic pressure, with shear and bulk viscosity,
with energy flowing between fluid elements, and with a frozen-in
magnetic field.  The stress-energy tensor for such a fluid is
%
\begin{align}
\tag{2.5.21}
T^{\alpha\beta}
&= \rho u^\alpha u^\beta + (p - \zeta\theta)\, h^{\alpha\beta}
   - 2\eta\, \sigma^{\alpha\beta}
   + q^\alpha u^\beta + u^\alpha q^\beta + \pi^{\alpha\beta}
\notag \\
&\quad
   + \frac{1}{8\pi}(B^2 u^\alpha u^\beta + B^2 h^{\alpha\beta}
   - 2B^\alpha B^\beta).
\end{align}

\noindent
The term $\rho u^\alpha u^\beta$ is the total density of mass-energy
(excluding only that of the frozen-in B-field) as measured in the LRF.
The term $p\, h^{\alpha\beta}$ is the isotropic pressure that would be
measured in the LRF if the gas were not changing volume (if $\theta$
were zero).  The quantities $\zeta$ and $\eta$ are the
\emph{coefficients of bulk viscosity and of dynamic viscosity},
respectively.  The term $-\zeta\theta\, h^{\alpha\beta}$ is the
isotropic viscous stress which resists isotropic expansion
($\theta > 0$) or compression ($\theta < 0$) of the fluid.
The term $-2\eta\, \sigma^{\alpha\beta}$ is the viscous shear stress
which resists shearing motions.  The term
$q^\alpha u^\beta + u^\alpha q^\beta$ is the energy density and
stresses associated with the flowing energy $\mathbf{q}$, as measured
in the LRF, are here neglected by comparison with $\rho u^\alpha u^\beta$
and $p\, h^{\alpha\beta}$.

(\emph{Note:} the energy density and stresses associated with the
flowing energy $\mathbf{q}$, as measured in the LRF, are here
neglected by comparison with $\rho u^\alpha u^\beta$ and
$p\, h^{\alpha\beta}$.  This neglect is valid with enormous accuracy in
most contexts of interest in these lectures.  An exception is the
radiation pressure which produces ``self-regulation'' of accretion onto
black holes when the inflowing mass is sufficiently large; see
\S\S\,4.5 and 5.13.)  The term
%
\begin{equation}
\tag{2.5.22}
T^{\alpha\beta}_{\text{MAG}}
= \frac{1}{8\pi}(B^2 u^\alpha u^\beta + B^2 h^{\alpha\beta}
  - 2B^\alpha B^\beta)
\end{equation}
%
is the Maxwell stress-energy associated with the frozen-in B-field (in
LRF: energy density $B^2/8\pi$, pressure $B^2/8\pi$ orthogonal to
field lines; tension $-B^2/8\pi$ along field lines).

\bigskip
\noindent
\emph{Equations of hydrodynamics}

\smallskip
\noindent
The fundamental equations governing the motion of a fluid in a given
gravitational field (spacetime geometry) are (i) the law of
\emph{baryon conservation}
%
\begin{equation}
\tag{2.5.23}
\nabla \cdot (n\,\mathbf{u}) = 0,
\quad\text{i.e.}\quad dn/d\tau = -\theta\, n,
\end{equation}
%
or equivalently rest-mass conservation
%
\begin{equation}
\tag{2.5.23$'$}
\nabla \cdot (\rho_0\, \mathbf{u}) = 0,
\quad\text{i.e.}\quad d\rho_0/d\tau = -\theta\, \rho_0\,;
\end{equation}

\noindent
(ii) the law of \emph{local energy conservation}
%
\begin{equation}
\tag{2.5.24}
\mathbf{u}\cdot(\nabla\cdot\mathbf{T}) = 0;
\end{equation}

\noindent
(iii) the \emph{Euler equations} (i.e.\ law of local momentum
conservation)
%
\begin{equation}
\tag{2.5.25}
\mathbf{h}\cdot(\nabla\cdot\mathbf{T}) = 0;
\end{equation}

\noindent
(iv) the laws of thermodynamics, (2.5.11)--(2.5.16); (v) the laws of
energy transport for the frozen-in B-field, (2.5.20); the laws of
energy transport, which govern $\mathbf{q}$ (see \S\,2.6, 2.6.43).

\bigskip
\noindent
\emph{Law of local energy conservation}

\smallskip
\noindent
When one evaluates the law of local energy conservation,
$\mathbf{u}\cdot(\nabla\cdot\mathbf{T}) = 0$,
for the stress-energy tensor (2.5.21), one finds that the Maxwell
stress-energy
%
\begin{equation}
\tag{2.5.26}
\mathbf{u}\cdot(\nabla\cdot\mathbf{T}_{\text{MAG}}) = 0
\end{equation}
%
gives zero contribution (work done to compress magnetic field is
precisely equal to increase in magnetic field energy); and that the
remainder of the stress-energy tensor gives
%
\begin{equation}
\tag{2.5.27}
d\rho/d\tau = -(\rho + p)\,\theta
+ \zeta\,\theta^2 + 2\eta\,\sigma_{\alpha\beta}\,\sigma^{\alpha\beta}
- \nabla\cdot\mathbf{q} - \mathbf{a}\cdot\mathbf{q}\,.
\end{equation}

\noindent
Here $d/d\tau$ is derivative with respect to proper time along the
world lines of the fluid; i.e., in the language of the differential
geometer, $d/d\tau = \mathbf{u}$.  The terms $-(\rho + p)\,\theta$ is
the increase in mass-energy density due to compression.  The term
$\zeta\,\theta^2$ and $2\eta\,\sigma_{\alpha\beta}\,\sigma^{\alpha\beta}$
are the increases in mass-energy density due to viscous heating
(conversion of relative kinetic energy of adjacent fluid elements into
heat).  The term $-\nabla\cdot\mathbf{q}$ is the influx of
mass-energy from neighboring fluid elements.  The term
$-\mathbf{a}\cdot\mathbf{q}$ is a special relativistic correction to
the ``redshift'' of $\mathbf{q}$.  To understand this term, consider a
flux of energy $\mathbf{q}$---or, equivalently, with the inertia of
the energy flux $\mathbf{q}$ associated with it, a flux of mass-energy
from the viewpoint of an accelerated fluid (photons do not interact with
them).  The acceleration gives rise to a redshift (a gravitational
redshift) and hence to a nonzero $\nabla\cdot\mathbf{q}$ (no
interaction between fluid and photons), one must include the correction
factor $-\mathbf{a}\cdot\mathbf{q}$.  For a similar reason, a similar
relativistic correction appears in the law of heat conduction.

\bigskip
\noindent
\emph{Euler equation for a perfect fluid}

\smallskip
\noindent
Consider a perfect fluid, flowing adiabatically through spacetime
($\zeta = \eta = \mathbf{B} = \mathbf{q} = 0$).  In this case the
Euler equations $\mathbf{h}\cdot(\nabla\cdot\mathbf{T}) = 0$ reduce to
%
\begin{equation}
\tag{2.5.30}
(\rho + p)\,\mathbf{a} = -\nabla p - \mathbf{u}\,(\mathbf{u}\cdot\nabla p).
\end{equation}
%
In words:
\[
(\rho + p)\,\mathbf{a}
= -\left(\!\text{inertial mass per}\atop\text{unit volume}\!\right)
  \times\text{(4-acceleration)}
= -\left(\!\text{pressure gradient}\atop\text{projected orthogonal to }\mathbf{u}\!\right).
\]
